% ĐỀ THI TỐT NGHIỆP THPT (kiểu Bộ GD&ĐT): ba phần trắc nghiệm, không tự luận. In đen trắng.
% \dethamkhao đổi thành ĐỀ THAM KHẢO.
\documentclass[12pt]{article}
\usepackage{thptqg}
\begin{document}
\namthi{2026}
\dechinhthuc
\made{0101}
\taodaude

\phantn
\begin{tn}
Tập xác định của hàm số $y=\tan x$ là
\paa{$\mathbb{R}$.}{\dung{$\mathbb{R}\setminus\left\{\frac{\pi}{2}+k\pi \mid k\in\mathbb{Z}\right\}$.}}{$\mathbb{R}\setminus\{k\pi \mid k\in\mathbb{Z}\}$.}{$[-1;1]$.}
\end{tn}

\phands
\begin{ds}
Cho hàm số $y=x^3-3x$.
\ydung{Hàm số có hai điểm cực trị.}
\ysai{Hàm số đồng biến trên $\mathbb{R}$.}
\ydung{$y(2)=2$.}
\ysai{Đồ thị cắt trục $Ox$ tại một điểm.}
\end{ds}

\phantln
\begin{tln}
Cho $\int_0^1 f(x)\,dx=2$. Tính $\int_0^1 3f(x)\,dx$.
\traloi{6}
\end{tln}
\bangdapan
\hetde
\end{document}
